\section{Experimental Setting}



\subsection{Tasks}
We adopt classification tasks from the benchmark GLUE \cite{wang2018glue} and those from MBPA++ \cite{huang2021continual,de2019episodic}. 
We select dissimilar tasks to form the task sequences, i.e. there are no overlap labels between each task.   The tasks contain 1) CoLA \cite{warstadt2019neural}, requiring the model to  determine whether a sentence is linguistically acceptable; 2) MNLI \cite{williams2017broad} containing 433k sentence pairs annotated with textual entailment information; 3) QNLI\footnote{\href{https://quoradata.quora.com/First-Quora-Dataset-Release-Question-Pairs}{https://quoradata.quora.com/First-Quora-Dataset-Release-Question-Pairs}}, requiring deciding whether the \textit{answer} answers the \textit{question}; 4) QQP, (parsed from SQuAD \cite{rajpurkarsquad}), testing whether a pair of Quora questions are synonymous; 5) Yelp \cite{zhang2015character}, requiring detecting the sentiment of a sentence; 6) AG \cite{zhang2015character}, requiring to classify the topics of the news. 

The sequences can be divided into 2 types with respect to the task lengths: 1) a sequence of 4 classification tasks containing AG, Yelp, QNLI, and MRPC; 2) a sequence of 6 classification tasks containing AG, MRPC, MNLI, CoLA, Yelp, and QNLI.  Without losing generality the orders are randomly selected and the task orders for experiments are shown in Table \ref{seq}.


\begin{table}[] \small
\centering
\begin{tabular}{ll}
\hline
\multicolumn{2}{l}{Orders}         \\ \hline

1      & AG $\rightarrow$ Yelp $\rightarrow$ QNLI$\rightarrow$ MRPC \\
2      & MRPC $\rightarrow$ QNLI $\rightarrow$ Yelp $\rightarrow$AG  \\
3      & QNLI $\rightarrow$Yelp $\rightarrow$MRPC$\rightarrow$AG \\ 
4 & AG$\rightarrow$MRPC $\rightarrow$CoLA$\rightarrow$MNLI $\rightarrow$Yelp$\rightarrow$ QNLI \\
5 & QNLI $\rightarrow$Yelp $\rightarrow$MNLI$\rightarrow$CoLA $\rightarrow$MRPC$\rightarrow$ AG \\
6 & MNLI $\rightarrow$ AG $\rightarrow$QNLI$\rightarrow$ 
 MRPC $\rightarrow$Yelp$\rightarrow$ CoLA \\
\hline
\end{tabular}
\caption{Different  task orders for our experiments.}
\label{seq}
    \vspace{-5mm}
\end{table}

\begin{table*}[]\small
\centering
\begin{tabular}{l|cc|cc|cc|cc|cc|cc}
\hline 
Model  & \multicolumn{2}{c|}{Order 1} & \multicolumn{2}{c|}{Order 2} & \multicolumn{2}{c|}{Order 3} & \multicolumn{2}{c|}{Order 4} & \multicolumn{2}{c}{Order 5}& \multicolumn{2}{c}{Order 6} \\ \hline
       & ACC           & BWT         & ACC          & BWT         & ACC          & BWT         & ACC          & BWT         & ACC          & BWT   & ACC          & BWT       \\
       \hline
Joint & 83.09 & -  &  83.09 & - & 83.09 & -  & 83.06 & - & 83.06 & - & 83.06 & -  \\
\cdashline{1-13}[2pt/4pt]
CE     & 54.74 & -36.96 & 59.43 & -31.11  & 51.73  &  -41.22 & 56.05  & -29.71 & 49.70 & -32.32 & 48.64 &-39.96 \\
CL     &  68.10  & -14.74    & 59.95   & -23.13     &  60.13    &  -22.72    &  62.35   &-24.28   & 58.65    & -27.26  & 63.36  &  -18.7           \\
\cdashline{1-13}[2pt/4pt]
% \hline
% \cdashline{1-13}[2pt/4pt]
CE-MAS &  59.91    & -23.53   &  61.21
    &  -24.45  &  61.75  & -19.25            &  62.52   & -19.25    & 54.77          &  -17.26  & 58.30 & -32.58
         \\
CL-MAS & 68.99   &   -1.37     &  71.83            &  -3.55    &  73.61   &  -6.94        & 69.95           & -2.89  & 69.27    & -6.27   & 68.93  & -2.44         \\ 
CE-EWC & 66.11      & -22.83     &  71.44            & -14.56     & 69.66    & -16.47            &  62.95            & -22.42    & 59.22     & -25.79  &61.87  &-22.81             \\
CL-EWC &  73.19      &   -0.59     &  75.89            &    -2.04         &  73.52            &   -5.97      &  66.74    & -3.65      & 68.83      &  -6.27  & 68.56  &  -4.00         \\
CE-LwF &  72.09     & -13.26      & 72.13             &  -14.39       & 73.54      &  -12.36       &  68.23     & -13.84     &  63.15    & -22.72  & 67.92  & -17.38            \\
CL-LwF &  76.53     &  -0.33      & 79.15      &  -3.71           &  79.58            &   -3.23          & 68.24            & -5.34   &    72.39     & -9.83    & 71.48  & -5.97        \\
% \cdashline{1-13}[2pt/4pt]

\cdashline{1-13}[2pt/4pt]
CE-ER  &  76.83     & -9.13   & 76.60             & -10.39   & 76.90   &  -12.59           &  75.08    &  -10.05   & {76.51}        & -6.39   & 76.15  & -14.61            \\
IDBR& 75.70  & -3.16  & 73.62  & -7.43  &  75.11  & -3.65 & 65.40 & -10.56 & 69.94 & -5.96  & 66.30  & -10.86\\
% DERPP&&&&&&&&&&&\\
Co$^2$L&70.58 & -2.07 & 74.02 & -7.52 & 74.10 & -7.68 & 64.31 & -3.57 & 65.04 & -10.62 & 64.94 &-15.65 \\
\cdashline{1-13}[2pt/4pt]


% CL-ER  &  79.02       &  -4.89  &   80.34           &  -3.85   & 79.95   &  -4.12      &  77.93            &  -5.17    &  76.91  &  -5.51   & 74.83           \\
% CL-IRD & \textbf{77.74}  & -5.23  & \textbf{80.19}  & -2.97  &  \textbf{80.12}   &  -2.85   & \textbf{77.80}  & -4.49  & {76.51}   &  -4.72   & \textbf{74.73}   & -4.32 \\   
% \textit{\textbf{Ours}}  &&&&&&&&&&&\\
\textbf{SCCL} & \textbf{79.20} & -2.93 & \textbf{80.05}  & -3.07    &  \textbf{80.24}  &  -3.51  &  \textbf{78.36}   &  0.57   & \textbf{79.00}   & -3.39  & \textbf{78.55}  & -3.75  \\
\ \ {{w/o MR}} & 77.19 & -5.34  & 78.63  &  -4.59   &  80.27  &  -2.91   &  75.65   & -2.91     &  71.22  &  -11.62  &  74.64   & -7.30             \\
\ \ {{w/o IRD}} & 77.57 & -5.33 & 79.73  & -2.48    &  79.48  &  -3.33  &  73.87   &  -6.62   & 76.89   & -4.10  & 74.32  & -4.03   \\
\hline
\end{tabular}
\caption{Continual Learning results on 6 different tasks.  `CE' refers to the standard cross-entropy-based methods, and `CL' refers to extended  contrastive-learning-based methods with continual learning strategies. `-' for not acquirable. All the results are averaged on 5 different random seeds. }
\label{results}
    \vspace{-5mm}
\end{table*}

\subsection{Evaluation Metrics}
We adopt the metrics of average accuracy (ACC) and backward transfer (BWT) to evaluate the performance of the continual learning model \cite{lopez2017gradient}. The model trained after the task $T^i$ is evaluated on the test set of earlier tasks $T^j$ ($j<=i$), and the test accuracy is denoted as $R^i_j$. The metrics are shown as follows:
% \vspace{-2mm}
\begin{equation}
    ACC = \frac{1}{n} \sum_{i=1}^{n} R^n_i
\end{equation}

\begin{equation}
    BWT = \frac{1}{n-1} \sum_{i=1}^{n-1} R^n_i - R^i_i,
\end{equation}
where the former evaluates the overall performance of the final trained model, and the latter calculates the knowledge forgetting during the continual training procedure.


% \noindent \textbf{\textrm{RoBERTa}} pre-trains the model BERT on over 160GB English corpora by adopting improved techniques such as dynamic sampling, large mini-batches, and masked language modeling without the task of next sentence prediction.



 \subsection{Baselines} We not only compare our model with several CE-based continual learning methods but extend training strategies of them to our contrastive learning framework (i.e. training with contrastive learning and inferring with kNN) to verify the effectiveness of contrastive learning in mitigating CF. We also compare our model with the competitive models IRDB \cite{huang2021continual} and Co$^2$L \cite{cha2021co2l}.  The shared hyper-parameters are kept the same as SCCL in baselines. The model details are as follows:

\begin{itemize}
     \vspace{-2mm}
     \item \noindent \textbf{Fine-tune} (CE, CL) \cite{yogatama2019learning} modifies the parameters of the pre-trained language model to adapt to a new task without any augmented strategies and additional loss.
     \vspace{-2mm}
     \item \noindent \textbf{Experience Replay} (ER) \cite{riemer2019learning} stores a small subset of samples from previous tasks and replays those to prevent models from forgetting past knowledge. 
    \vspace{-2mm}
     \item \noindent \textbf{Elastic Weight Consolidation} (EWC) \cite{kirkpatrick2017overcoming} slows down the updates of the optimal parameters for previous tasks by extending the loss function with a regularization term. 
 \vspace{-2mm} 
     \item \noindent \textbf{Memory Aware Synapses} (MAS) \cite{aljundi2018memory} slows down the update according to the importance weight of each parameter in the network, i.e. the sensitivity of the output function to a parameter change. 
\vspace{-2mm} 
     \item \noindent \textbf{Learning Without Forgetting} (LwF) \cite{li2017learning,9156964}  aims to keep the model output of current data close to those of the previous model. 
  \vspace{-2mm} 
     \item \noindent \textbf{IDBR} \cite{huang2021continual} uses information  disentanglement regularization to encode task-specific information and general information individually, which are jointly considered for classification.
 \vspace{-2mm} 
      \item \noindent \textbf{Co$^2$L} \cite{cha2021co2l} uses an asymmetric supervised contrastive learning method to learn representations and trains  a decoupled  layer for inference.
  \vspace{-2mm}
      \item \noindent \textbf{Multi-task Training} (Joint) \cite{NEURIPS2020_3fe78a8a}  trains on all the tasks simultaneously, i.e. the data of different tasks are mixed up for training. It does not suffer from catastrophic forgetting and represents an upper bound on model performance. 
\end{itemize}










 % \begin{enumerate}
 % \vspace{-1mm}
 %     \item \noindent \textbf{Fine-tune} \cite{yogatama2019learning} modifies the parameters of the pre-trained language model to adapt to a new task without any augmented strategies and additional loss.
 %     \vspace{-1mm}
     
 % \end{enumerate}




% We follow the modified version proposed by \citet{9156964}, regularizing the output embeddings of current data  with respect to their embeddings in the previous task.
 % \begin{equation}
 %     \mathcal{L}_{LwF} = ||z_i^{m} - z_i^{m-1}||, 
 % \end{equation}
 % where $||\cdot||$ refers to the Frobenius norm. $z_i^{m}$ and $z_i^{m-1}$ refer to representations obtained from the current model and previous model, respectively.


% The extended loss function is as follows:
% \begin{equation}
%      \mathcal{L}_{EWC} = \sum_p \frac{1}{2} \mathcal{F}_p^{m-1}(\theta^m_p - \theta^{m-1}_p)^2, 
%  \end{equation}
% where $\mathcal{F}^{m-1}$ is the Fisher information matrix calculated after training on the previous task $t-1$, and $\theta_p$ is the sum of the model parameters.

% The function is as follows:
% \begin{equation}
% \mathcal{L}_{MAS} = \sum_p \frac{1}{2} \Delta_p(\theta^m_p - \theta^{m-1}_p)^2, 
% \end{equation}
% where $\Delta_p$ is estimated through the sensitivity of the squared $l_2$ norm of the predicted function output.


\subsection{Implementation Details}
 We adopt the officially released \textit{roberta-base} from HuggingFace \footnote{\href{https://huggingface.co/}{https://huggingface.co/}} as our backbone network.  We train our model on 1 GPU (A100 80G) using the Adam optimizer \cite{kingma2014adam}. For all the models, the batch size is 96,  the learning rate is 3e-5, and the scheduler is set linear. We train our model 10 epochs for each task. Following \citet{huang2021continual}, we select 4,000 samples for each label in training. The hyper-parameters of temperatures $\kappa$ is 0.2, $\tau^*$ is 0.2, and $T$ is 5, and the number of nearest neighbors $k$ is 10. The  memory size for each task is set to 200 (2.5\% of the training data) and the memory replay frequency $f$ is 100.  Through the training of our model, no development set is applied to find the best checkpoints, but stop until the training step is reached.






